% ============================================================
%  glossary.tex — Glosario de Términos (65 entradas)
%  Generado automáticamente desde el documento original.
%  Para imprimir todas las entradas usamos \glsaddall en main.tex
% ============================================================

\newglossaryentry{AAD}{name={AAD}, description={Additional Authenticated Data (Datos Adicionales Autenticados). Información que no se cifra pero cuya integridad se protege mediante la etiqueta de autenticación del cifrado autenticado GCM.}}
\newglossaryentry{AAL}{name={AAL}, description={Authentication Assurance Level (Nivel de Aseguramiento de Autenticación). Clasificación del NIST SP 800-63B que define tres niveles de confianza en la autenticación: AAL1 (factor único), AAL2 (dos factores) y AAL3 (dos factores con hardware criptográfico).}}
\newglossaryentry{ACR2SE}{name={ACR2-SE}, description={Autoridad Certificadora Raíz Segunda de la Secretaría de Economía. Raíz de la jerarquía PKI mexicana bajo la cual operan los PSC autorizados.}}
\newglossaryentry{AEAD}{name={AEAD}, description={Authenticated Encryption with Associated Data (Cifrado Autenticado con Datos Asociados). Esquema criptográfico que provee simultáneamente confidencialidad, integridad y autenticación del mensaje. AES-GCM es una instancia de AEAD.}}
\newglossaryentry{AES}{name={AES}, description={Advanced Encryption Standard (Estándar de Cifrado Avanzado). Algoritmo de cifrado simétrico por bloques estandarizado en FIPS 197. AES-256 utiliza claves de 256 bits.}}
\newglossaryentry{API}{name={API}, description={Application Programming Interface (Interfaz de Programación de Aplicaciones). Conjunto de definiciones y protocolos que permite la comunicación entre componentes de software.}}
\newglossaryentry{Argon2id}{name={Argon2id}, description={Variante híbrida de la función de hashing de contraseñas Argon2, ganadora de la Password Hashing Competition (2015) y estandarizada en RFC 9106. Combina resistencia a ataques de canal lateral (Argon2i) con resistencia a hardware especializado (Argon2d).}}
\newglossaryentry{ASIC}{name={ASIC}, description={Application-Specific Integrated Circuit (Circuito Integrado de Aplicación Específica). Hardware diseñado para ejecutar una tarea específica, utilizado en ataques de fuerza bruta contra funciones hash.}}
\newglossaryentry{ASN1}{name={ASN.1}, description={Abstract Syntax Notation One (Notación de Sintaxis Abstracta Uno). Estándar internacional para definir estructuras de datos criptográficas, utilizado en certificados X.509 y sellos de tiempo RFC 3161.}}
\newglossaryentry{backend}{name={backend}, description={Capa lógica del servidor que procesa la lógica de negocio, el acceso a datos y los servicios criptográficos, en contraposición al frontend que interactúa con el usuario.}}
\newglossaryentry{bcrypt}{name={bcrypt}, description={Función de hashing de contraseñas diseñada en 1999, basada en el cifrador Blowfish. Incorpora un factor de costo ajustable pero carece de dureza de memoria configurable.}}
\newglossaryentry{CA}{name={CA}, description={Certificate Authority (Autoridad Certificadora). Entidad de confianza responsable de emitir, renovar y revocar certificados digitales dentro de una infraestructura de clave pública (PKI).}}
\newglossaryentry{CAdEST}{name={CAdES-T}, description={CMS Advanced Electronic Signatures with Timestamp. Formato de firma electrónica avanzada que incorpora un sello de tiempo dentro de una estructura CMS (Cryptographic Message Syntax).}}
\newglossaryentry{CRL}{name={CRL}, description={Certificate Revocation List (Lista de Certificados Revocados). Documento firmado por una CA que enumera los certificados revocados antes de su fecha de expiración, conforme a RFC 5280.}}
\newglossaryentry{CSR}{name={CSR}, description={Certificate Signing Request (Solicitud de Firma de Certificado). Mensaje enviado a una CA para solicitar la emisión de un certificado digital, conteniendo la clave pública del solicitante.}}
\newglossaryentry{CSPRNG}{name={CSPRNG}, description={Cryptographically Secure Pseudo-Random Number Generator (Generador de Números Pseudoaleatorios Criptográficamente Seguro). Generador cuya salida es computacionalmente indistinguible de datos aleatorios verdaderos.}}
\newglossaryentry{DEK}{name={DEK}, description={Data Encryption Key (Clave de Cifrado de Datos). En un esquema de envelope encryption, clave simétrica generada aleatoriamente para cifrar un documento individual.}}
\newglossaryentry{DER}{name={DER}, description={Distinguished Encoding Rules. Reglas de codificación binaria para estructuras ASN.1, utilizadas en certificados X.509 y sellos de tiempo.}}
\newglossaryentry{DOF}{name={DOF}, description={Diario Oficial de la Federación. Órgano del Gobierno de México para la publicación de leyes, reglamentos y acuerdos oficiales.}}
\newglossaryentry{ECDSA}{name={ECDSA}, description={Elliptic Curve Digital Signature Algorithm (Algoritmo de Firma Digital de Curva Elíptica). Esquema de firma basado en la dificultad del logaritmo discreto sobre curvas elípticas.}}
\newglossaryentry{efirma}{name={e.firma}, description={Certificado de firma electrónica avanzada emitido por el Servicio de Administración Tributaria (SAT) para personas físicas y morales en México. Utiliza RSA-2048 con SHA-256.}}
\newglossaryentry{endpoint}{name={endpoint}, description={Punto de acceso de una API, identificado por una URL y un método HTTP (por ejemplo, GET /pki/crl.pem).}}
\newglossaryentry{envelopeencryption}{name={envelope encryption}, description={Cifrado de envolvente. Esquema de gestión de claves donde una clave maestra (KEK) cifra claves de datos individuales (DEK), cada una asociada a un documento específico.}}
\newglossaryentry{FIPS}{name={FIPS}, description={Federal Information Processing Standards (Estándares Federales de Procesamiento de Información). Conjunto de estándares del gobierno de Estados Unidos para sistemas informáticos. FIPS 180-4 define SHA-256; FIPS 197 define AES.}}
\newglossaryentry{framework}{name={framework}, description={Marco de trabajo o plataforma de desarrollo que proporciona una estructura predefinida para construir aplicaciones de software.}}
\newglossaryentry{GCM}{name={GCM}, description={Galois/Counter Mode (Modo Galois/Contador). Modo de operación de cifrado por bloques que provee cifrado autenticado (AEAD), especificado en NIST SP 800-38D.}}
\newglossaryentry{GPU}{name={GPU}, description={Graphics Processing Unit (Unidad de Procesamiento Gráfico). Procesador diseñado para cómputo paralelo masivo, utilizado también en ataques de fuerza bruta contra funciones hash.}}
\newglossaryentry{hash}{name={hash}, description={Resumen criptográfico o huella digital. Valor de longitud fija producido por una función hash criptográfica a partir de una entrada de longitud arbitraria. Cualquier modificación en la entrada produce un hash completamente diferente.}}
\newglossaryentry{HKDF}{name={HKDF}, description={HMAC-based Key Derivation Function (Función de Derivación de Claves basada en HMAC). Mecanismo para derivar múltiples claves criptográficas a partir de una clave maestra y un contexto único.}}
\newglossaryentry{HMAC}{name={HMAC}, description={Hash-based Message Authentication Code (Código de Autenticación de Mensajes basado en Hash). Construcción criptográfica que combina una función hash con una clave secreta para verificar integridad y autenticidad.}}
\newglossaryentry{HOTP}{name={HOTP}, description={HMAC-Based One-Time Password (Contraseña de Un Solo Uso basada en HMAC). Algoritmo definido en RFC 4226 que genera códigos de un solo uso a partir de un secreto compartido y un contador.}}
\newglossaryentry{IETF}{name={IETF}, description={Internet Engineering Task Force. Organismo internacional que desarrolla y publica estándares de protocolos de Internet, incluyendo las RFC.}}
\newglossaryentry{IV}{name={IV}, description={Initialization Vector (Vector de Inicialización). Valor aleatorio o único utilizado junto con una clave para iniciar un algoritmo de cifrado, garantizando que el mismo texto plano produzca diferentes textos cifrados.}}
\newglossaryentry{JWT}{name={JWT}, description={JSON Web Token. Estándar abierto (RFC 7519) para la creación de tokens de acceso firmados digitalmente en formato JSON, utilizados para autenticación y autorización.}}
\newglossaryentry{KEK}{name={KEK}, description={Key Encryption Key (Clave de Cifrado de Clave). En un esquema de envelope encryption, clave maestra utilizada exclusivamente para cifrar y descifrar las DEKs.}}
\newglossaryentry{legacyuse}{name={legacy-use}, description={Uso heredado. Clasificación del NIST para algoritmos o tamaños de clave que son aceptables para procesar información existente pero no recomendados para proteger información nueva.}}
\newglossaryentry{memoryhardness}{name={memory-hardness}, description={Dureza de memoria. Propiedad de una función criptográfica que requiere una cantidad configurable de memoria durante su ejecución, encareciendo proporcionalmente los ataques con hardware paralelo.}}
\newglossaryentry{MFA}{name={MFA}, description={Multi-Factor Authentication (Autenticación por Múltiples Factores). Mecanismo que requiere la presentación de al menos dos categorías de evidencia de identidad: algo que el usuario sabe, posee o es.}}
\newglossaryentry{middleware}{name={middleware}, description={Capa de software intermedia que intercepta y procesa peticiones HTTP entre el cliente y la lógica de la aplicación, realizando funciones transversales como autenticación, registro de actividad o limitación de frecuencia.}}
\newglossaryentry{NIST}{name={NIST}, description={National Institute of Standards and Technology (Instituto Nacional de Estándares y Tecnología). Organismo del gobierno de Estados Unidos que publica estándares criptográficos y guías de seguridad.}}
\newglossaryentry{NOM}{name={NOM}, description={Norma Oficial Mexicana. Regulación técnica de observancia obligatoria emitida por dependencias del gobierno federal de México. NOM-151-SCFI-2016 establece los requisitos para la conservación de mensajes de datos.}}
\newglossaryentry{nonce}{name={nonce}, description={Número de uso único (del inglés number used once). Valor que no debe repetirse con la misma clave en operaciones criptográficas. En AES-GCM, su reutilización compromete catastróficamente el esquema.}}
\newglossaryentry{OCSP}{name={OCSP}, description={Online Certificate Status Protocol (Protocolo de Estado de Certificado en Línea). Protocolo que permite consultar en tiempo real el estado de revocación de un certificado ante un servidor autorizado.}}
\newglossaryentry{OID}{name={OID}, description={Object Identifier (Identificador de Objeto). Secuencia numérica única que designa un algoritmo, atributo u objeto en el estándar ASN.1. Ejemplo: 1.2.840.113549.1.1.11 identifica sha256WithRSAEncryption.}}
\newglossaryentry{OWASP}{name={OWASP}, description={Open Worldwide Application Security Project. Proyecto abierto de referencia internacional en seguridad de aplicaciones web, que publica guías de mejores prácticas y listas de vulnerabilidades.}}
\newglossaryentry{PAdEST}{name={PAdES-T}, description={PDF Advanced Electronic Signatures with Timestamp. Formato de firma electrónica avanzada embebida dentro de un documento PDF, con sello de tiempo incorporado.}}
\newglossaryentry{PEM}{name={PEM}, description={Privacy-Enhanced Mail. Formato de codificación Base64 para certificados, claves criptográficas y listas de revocación, delimitado por encabezados como -----BEGIN CERTIFICATE-----.}}
\newglossaryentry{PHC}{name={PHC}, description={Password Hashing Competition (Competencia de Hashing de Contraseñas). Concurso criptográfico internacional celebrado entre 2013 y 2015 con 24 candidatos, cuyo ganador fue Argon2.}}
\newglossaryentry{PKI}{name={PKI}, description={Public Key Infrastructure (Infraestructura de Clave Pública). Conjunto de entidades, políticas y procedimientos para crear, gestionar, distribuir y revocar certificados digitales.}}
\newglossaryentry{PKCS1}{name={PKCS\#1}, description={Public-Key Cryptography Standards \#1. Estándar que define el formato y los esquemas de firma y cifrado RSA, incluyendo los esquemas de relleno (padding) v1.5 y PSS.}}
\newglossaryentry{PSC}{name={PSC}, description={Prestador de Servicios de Certificación. Entidad autorizada por la Secretaría de Economía de México para emitir certificados digitales, sellos de tiempo y constancias de conservación bajo la jerarquía ACR2-SE.}}
\newglossaryentry{ratelimiting}{name={rate limiting}, description={Limitación de frecuencia. Mecanismo de seguridad que restringe el número de peticiones o intentos que un usuario puede realizar en un período de tiempo, previniendo ataques de fuerza bruta.}}
\newglossaryentry{RBAC}{name={RBAC}, description={Role-Based Access Control (Control de Acceso Basado en Roles). Modelo de autorización que asigna permisos a roles predefinidos y asocia cada usuario a uno o más roles.}}
\newglossaryentry{REST}{name={REST}, description={Representational State Transfer (Transferencia de Estado Representacional). Estilo arquitectónico para diseñar servicios web basado en recursos identificados por URLs y operados mediante métodos HTTP estándar.}}
\newglossaryentry{RFC}{name={RFC}, description={Request for Comments (Solicitud de Comentarios). Serie de documentos técnicos publicados por el IETF que definen estándares y protocolos de Internet. Cada RFC se identifica con un número único (por ejemplo, RFC 3161).}}
\newglossaryentry{RSA}{name={RSA}, description={Rivest-Shamir-Adleman. Algoritmo de criptografía asimétrica basado en la dificultad de factorizar números enteros grandes, utilizado para firma digital y cifrado de clave pública.}}
\newglossaryentry{salt}{name={salt}, description={Sal. Valor aleatorio único que se concatena con la contraseña antes de aplicar la función hash, garantizando que contraseñas idénticas produzcan hashes diferentes.}}
\newglossaryentry{sandbox}{name={sandbox}, description={Entorno de pruebas aislado que replica las funcionalidades de un sistema en producción sin afectar datos reales ni generar efectos legales.}}
\newglossaryentry{SAT}{name={SAT}, description={Servicio de Administración Tributaria. Órgano del gobierno de México responsable de la recaudación fiscal y emisor de la e.firma para personas físicas y morales.}}
\newglossaryentry{SHA256}{name={SHA-256}, description={Secure Hash Algorithm de 256 bits. Función hash criptográfica de la familia SHA-2 estandarizada en FIPS 180-4, que produce un resumen de 256 bits (32 bytes) y proporciona 128 bits de resistencia a colisiones.}}
\newglossaryentry{SP}{name={SP}, description={Special Publication (Publicación Especial). Serie de documentos técnicos del NIST que contienen guías, recomendaciones y estándares de seguridad. Ejemplo: NIST SP 800-57.}}
\newglossaryentry{TOTP}{name={TOTP}, description={Time-Based One-Time Password (Contraseña de Un Solo Uso basada en Tiempo). Algoritmo definido en RFC 6238 que genera códigos temporales a partir de un secreto compartido y la hora actual, utilizado como segundo factor de autenticación.}}
\newglossaryentry{TSA}{name={TSA}, description={Time-Stamp Authority (Autoridad de Sellado de Tiempo). Entidad de confianza que emite sellos de tiempo criptográficos conforme al protocolo RFC 3161, acreditando la existencia de un documento en un momento determinado.}}
\newglossaryentry{WebAuthn}{name={WebAuthn}, description={Web Authentication. Estándar del W3C y FIDO Alliance para autenticación sin contraseña basada en criptografía de clave pública, mediante llaves de seguridad físicas o biometría del dispositivo.}}
\newglossaryentry{X509}{name={X.509}, description={Estándar de la ITU-T que define el formato de los certificados digitales de clave pública, incluyendo la estructura de la cadena de confianza, las extensiones y los perfiles de listas de revocación (CRL).}}
